\documentclass[%
  10pt,% default; alternatives: 11pt, 12pt, 14pt
  %handout,% flatten overlays for the sake of printed handouts
]{../templates/pres_template/cognigronbeamer}

\usepackage{textcomp} % ensures TS1 symbols are available
\DeclareTextSymbolDefault{\textbullet}{TS1}

% for the frame header:
\facname{Faculty of Science\\ and Engineering} % optional
\group{Bio-inspired\\ Circuits and Systems} % optional
\inst{cognigron} % optional
\date{\today}
% \extralogo{...} % disabled in image-free version

% for the title frame:
\title{A Title}
\subtitle{RUG beamer-based presentations}
\author{A. Uthor}
\institute{University of Groningen\\ Faculty of Science and Engineering}

\begin{document}

% title slide
{
\setbeamertemplate{background canvas}{\titleback}
\begin{frame}%[plain]
\titlepage
\end{frame}
}

% section slide
\section{Big Menu}
%\label{sec:samples}

% standard slide
\begin{frame}
  \frametitle{Today's menu}
  Today we have a special menu.
\end{frame}

% itemize slide
\begin{frame}{Today's menu - first course}
  \begin{itemize}
    \item soup
      \begin{itemize}
      \item tortilla filled with meat and vegetables
      \item refried beans
      \end{itemize}
  \end{itemize}
\end{frame}

% Stepped enumerated list with pause
\begin{frame}{Stepped enumerated list with {\ttfamily pause}}
  \begin{enumerate}
    \item Is the soup good?\pause
    \item Is the main course good?\pause
    \item Are the tortillas good?
  \end{enumerate}
\end{frame}



% empty background; put inside braces to localize background redefinition
{
\setbeamertemplate{background canvas}{}
\begin{frame}[plain]{Empty background example}
  Empty background
\end{frame}
}



\end{document}